% =============================================================================
%  Chapters 3 & 4
%  Sentiment-Driven Stock Trend Prediction — Graduation Project
% =============================================================================

% NOTE: This file is meant to be \input{} or \include{} into the main .tex
%       document.  It assumes the preamble already loads:
%         amsmath, amssymb, graphicx, booktabs, longtable, hyperref,
%         algorithm2e (or algorithmicx), tikz, subcaption, listings, xcolor
%       and defines \PROJECT, \FNSPID, etc. macros as needed.

% ============================================================================
\chapter{Developed Approach and System Model}
\label{ch:approach}
% ============================================================================

This chapter presents the end-to-end system that was designed, implemented, and
iteratively refined over the course of this project.  We decompose the system
into three complementary models following standard systems-engineering practice:
a \emph{Data Model} describing what information flows through the system and how
it is represented (Section~\ref{sec:data-model}), a \emph{Structural Model}
describing the neural-network architectures that learn from that data
(Section~\ref{sec:structural-model}), and a \emph{Dynamic Model} describing the
temporal training and evaluation protocol that governs how the structural model
interacts with the data model over time (Section~\ref{sec:dynamic-model}).

% --------------------------------------------------------------------------
\section{Data Model}
\label{sec:data-model}
% --------------------------------------------------------------------------

The data model comprises four stages: (i)~raw data acquisition,
(ii)~sentiment extraction and aggregation, (iii)~technical indicator
computation, and (iv)~feature selection and caching.  Each stage is described
in full below.

% ·····································
\subsection{Raw Data Sources}
\label{sec:raw-data}
% ·····································

Two primary data sources feed the system:

\paragraph{Price Data.}
Daily Open--High--Low--Close--Volume (OHLCV) data for 200 S\&P\,500 constituent
stocks is retrieved from \texttt{yfinance}.  The date range spans
\textbf{2012-01-01 to 2019-12-31}, deliberately restricted to the era in which
the sentiment corpus has dense coverage.  A dedicated \texttt{PriceCache} module
(Section~\ref{sec:caching}) ensures that the 23\,GB download is performed only
once and subsequent runs hit a local pickle cache.

\paragraph{News Sentiment Corpus (FNSPID).}
The Financial News and Stock Price Integration Dataset
(FNSPID)~\cite{fnspid2024} provides roughly 29.7\,million financial news
headlines spanning 2009--2020.  Each record maps a headline to a ticker symbol
and a UTC timestamp.  We ingest this corpus through a streaming CSV reader that
filters by our 200-ticker universe and persists per-ticker Parquet files to
disk.

% ·····································
\subsection{Sentiment Extraction via FinBERT}
\label{sec:finbert}
% ·····································

Raw headlines carry no quantitative signal that a neural network can consume.
We therefore transform each headline into a three-dimensional probability vector
using \textbf{FinBERT}~\cite{finbert2019}, a domain-specific variant of the BERT
architecture~\cite{devlin2019bert}.

FinBERT was pre-trained on a large financial corpus (including Reuters TRC2 and
financial articles) and then fine-tuned on the Financial PhraseBank dataset for
three-class sentiment classification.  It inherits BERT's 12-layer
Transformer-encoder backbone with 110\,M parameters, but its vocabulary and
pre-training objective are tailored to financial language, making it
substantially more accurate than general-purpose sentiment classifiers on
financial text.

The extraction pipeline operates as follows:

\begin{enumerate}
  \item \textbf{Tokenisation.}  Each headline is tokenised using FinBERT's
        WordPiece tokeniser with a maximum sequence length of 512 sub-word
        tokens.  Shorter sequences are padded; longer ones are truncated.

  \item \textbf{Forward Pass.}  The tokenised input is passed through the
        12-layer Transformer encoder.  The \texttt{[CLS]} token's final hidden
        state (a 768-dimensional vector) is projected through a linear
        classification head to produce three raw logits.

  \item \textbf{Softmax Calibration.}  A softmax function converts the logits
        into a probability triple
        $(\hat{p}_{\text{pos}},\;\hat{p}_{\text{neg}},\;\hat{p}_{\text{neu}})$
        summing to one.

  \item \textbf{Sentiment Value.}  A scalar sentiment value is computed as the
        signed polarity:
        \begin{equation}
          s = \hat{p}_{\text{pos}} - \hat{p}_{\text{neg}}
          \label{eq:sentiment-value}
        \end{equation}
        This yields $s \in [-1, +1]$, where $s > 0$ indicates bullish sentiment
        and $s < 0$ indicates bearish sentiment.
\end{enumerate}

\noindent
Inference is performed in batches of 16 on a CUDA-enabled GPU with
\texttt{torch.no\_grad()} context to avoid gradient computation overhead.  The
resulting sentiment columns (label, score, positive, negative, neutral,
sentiment\_value) are persisted as per-ticker Parquet files for reuse.

% ·····································
\subsection{Daily Sentiment Aggregation}
\label{sec:aggregation}
% ·····································

Multiple headlines may appear for the same ticker on the same trading day.  The
\texttt{SentimentAggregator} module consolidates them into daily features via
one of three configurable strategies:

\begin{itemize}
  \item \textbf{Simple Average} (default for final experiments): computes the
        arithmetic mean, standard deviation, and count of all sentiment values
        $s$ falling on the same trading date.

  \item \textbf{Time-Weighted Average}: weights more recent intra-day headlines
        more heavily, assuming that later news supersedes earlier reports.

  \item \textbf{Sticky Sentiment with Exponential Decay}: implements a
        state-machine that maintains a running sentiment signal.  If the daily
        mean exceeds a significance threshold $|s| > 0.1$, the state is updated
        to the new value.  Otherwise, the previous state is decayed by a factor
        of $\alpha = 0.95$ per day.  This models the intuition that market
        sentiment persists beyond the publication date of a headline---traders do
        not forget yesterday's news instantaneously.
\end{itemize}

Additionally, \textbf{sentiment momentum} features are computed over rolling
windows of 3, 5, and 10 days, capturing the rate of change in the sentiment
signal.

\paragraph{No-News Day Handling.}
A critical design decision concerns trading days on which no headline exists for
a given ticker.  Rather than filling with zero (which would cause abrupt signal
drops), we apply an \textbf{exponential decay fill}: the last known sentiment
value is forward-filled and attenuated by $0.95^d$, where $d$ is the number of
consecutive no-news days.  This ensures that sentiment ``fades'' gradually rather
than vanishing, reflecting the real-world persistence of market narratives.

% ·····································
\subsection{Eight Candidate Sentiment Features}
\label{sec:sentiment-features}
% ·····································

From the aggregated daily sentiment, we engineer eight candidate features:

\begin{table}[htbp]
\centering
\caption{Eight candidate sentiment features derived from daily FinBERT scores.}
\label{tab:sentiment-features}
\begin{tabular}{@{}llp{6.5cm}@{}}
\toprule
\textbf{Feature}          & \textbf{Type}     & \textbf{Definition} \\
\midrule
\texttt{sent\_mean}       & Level             & Daily mean of $s$ \\
\texttt{sent\_dispersion} & Volatility        & Intra-day standard deviation of $s$ \\
\texttt{sent\_count\_log} & Activity          & $\log(1 + \text{news count})$ \\
\texttt{sent\_strength}   & Magnitude         & $|\texttt{sent\_mean}| \times \log(1+\text{count})$ \\
\texttt{sent\_surprise}   & Deviation         & $\texttt{sent\_mean} - \text{MA}_{20}(\texttt{sent\_mean})$ \\
\texttt{sent\_ema\_20d}   & Regime            & $\text{EMA}_{20}(\texttt{sent\_mean})$ \\
\texttt{sent\_momentum}   & Trend             & $\texttt{sent\_mean}_t - \texttt{sent\_mean}_{t-5}$ \\
\texttt{sent\_vol\_ratio} & Relative activity & $\texttt{count}_t\; /\; \text{MA}_{20}(\texttt{count})$ \\
\bottomrule
\end{tabular}
\end{table}

Not all eight features carry useful signal.  Feature selection
(Section~\ref{sec:feature-selection}) reduces this set to two per task
formulation.

% ·····································
\subsection{Technical Indicator Computation}
\label{sec:technical-indicators}
% ·····································

The technical feature pipeline evolved through seven major iterations.  The
final version (V7) uses the \texttt{pandas\_ta} library to compute approximately
100 indicators from raw OHLCV data, spanning the following categories:

\begin{itemize}
  \item \textbf{Trend Indicators}: Simple and Exponential Moving Averages (SMA, EMA),
        MACD and its histogram, Parabolic SAR, Aroon oscillator, ADX with
        directional indicators (DMP, DMN).

  \item \textbf{Momentum Indicators}: Relative Strength Index (RSI), Commodity
        Channel Index (CCI), Rate of Change (ROC), Stochastic oscillator, Chande
        Momentum Oscillator (CMO), TRIX.

  \item \textbf{Volatility Indicators}: Bollinger Bands (upper, middle, lower,
        bandwidth, percent-B), Average True Range (ATR), Donchian Channels,
        intraday range, gap percentage.

  \item \textbf{Volume Indicators}: On-Balance Volume (OBV), Accumulation/Distribution
        line, Chaikin Money Flow (CMF), Price--Volume Oscillator (PVO), volume
        ratio.

  \item \textbf{Domain-Knowledge Binary Features} (new in V7): Golden/Death Cross
        signals, RSI overbought/oversold zones, MACD crossover flags, Bollinger
        Band breakouts, volume confirmation signals.

  \item \textbf{Return Features}: 1-day, 5-day, and 10-day log returns, and
        derived features such as bear power and bull power.
\end{itemize}

% ·····································
\subsection{Feature Selection}
\label{sec:feature-selection}
% ·····································

From the ${\sim}100$ raw technical indicators, a two-stage selection process
reduces the feature set to a compact and informative subset.

\paragraph{Stage 1 --- SHAP-based Technical Feature Selection.}
A gradient-boosted tree (XGBoost) is trained on the full indicator set.
SHAP (SHapley Additive exPlanations) values~\cite{lundberg2017shap} are
computed to rank feature importance.  The top 20 technical features by mean
absolute SHAP value are retained:

\begin{quote}\small
\texttt{atr\_pct}, \texttt{intraday\_range}, \texttt{bb\_BBB\_5\_2.0\_2.0},
\texttt{gap}, \texttt{cci}, \texttt{high\_low\_pct}, \texttt{volume\_ratio},
\texttt{adx\_DMN\_14}, \texttt{return\_1d}, \texttt{aroon\_AROOND\_14},
\texttt{adx\_DMP\_14}, \texttt{obv}, \texttt{ad}, \texttt{bear\_power},
\texttt{roc\_10}, \texttt{pvo\_PVOh\_12\_26\_9}, \texttt{trix\_TRIXs\_30\_9},
\texttt{rsi\_14}, \texttt{aroon\_AROONU\_14}, \texttt{return\_5d}
\end{quote}

\paragraph{Stage 2 --- Sentiment Feature Selection via Permutation Importance.}
The eight candidate sentiment features (Table~\ref{tab:sentiment-features}) are
evaluated using RandomForest \textbf{permutation importance} rather than the
more commonly used Gini importance (Mean Decrease in Impurity).  This choice
is deliberate and motivated by three factors:

\begin{enumerate}
  \item \textbf{Bias toward high-cardinality features.}  Gini importance is
        computed during tree construction and is biased toward features with many
        unique values or high cardinality.  Continuous sentiment features would
        be systematically overvalued relative to binary or low-cardinality
        features, yielding misleading rankings.

  \item \textbf{Interaction effects.}  Gini importance measures how much a
        feature reduces impurity \emph{within the trained tree}, which is
        sensitive to the tree's specific split structure.  Permutation importance,
        by contrast, measures the \emph{out-of-sample} accuracy degradation when
        a feature's values are randomly shuffled, thereby capturing the feature's
        marginal contribution irrespective of how the model happened to use it.

  \item \textbf{Correlation robustness.}  When features are correlated (as
        several of our sentiment features are --- e.g.\ \texttt{sent\_mean} and
        \texttt{sent\_ema\_20d}), Gini importance distributes credit
        unpredictably among them.  Permutation importance exposes which features
        the model truly \emph{relies} on, because shuffling a redundant feature
        will not degrade performance if its correlated partner is intact.
\end{enumerate}

\noindent
The selection is task-specific:

\begin{itemize}
  \item \textbf{Classification task} (3-class trend, scoring = accuracy):
        \texttt{sent\_surprise} (\#1) and \texttt{sent\_ema\_20d} (\#2).
  \item \textbf{Regression task} (vol-adjusted return, scoring = $R^2$):
        \texttt{sent\_count\_log} (\#1) and \texttt{sent\_strength} (\#2).
\end{itemize}

The remaining six candidates are disabled to reduce noise injection.  The final
feature vector is therefore $20 + 2 = 22$ dimensional.

% ·····································
\subsection{Caching Architecture}
\label{sec:caching}
% ·····································

A multi-layer persistent caching system eliminates redundant computation across
experiment runs, effectively creating a \textbf{professional-grade quantitative
feature pool} that can be queried instantly:

\paragraph{Layer 1 --- Price Cache (\texttt{PriceCache}).}
The \texttt{PriceCache} class provides persistent pickle-based caching of raw
OHLCV data downloaded from \texttt{yfinance}.  Cache files are keyed by
\texttt{(start\_date, end\_date, suffix)} and stored under
\texttt{data/price\_cache/}.

\begin{itemize}
  \item On first invocation, \texttt{PriceCache.load()} returns \texttt{None}
        (cache miss), triggering a fresh download.
  \item \texttt{PriceCache.save()} serialises the resulting
        \texttt{Dict[str, DataFrame]} to a pickle file, where each key is a
        ticker symbol and each value is the full OHLCV DataFrame.
  \item On subsequent runs, \texttt{load()} deserialises the pickle and
        optionally filters to the requested ticker subset, avoiding the network
        round-trip entirely.
  \item An \texttt{overwrite} flag allows forced refresh when data staleness is
        suspected.
\end{itemize}

\paragraph{Layer 2 --- Sentiment Cache.}
FinBERT inference is computationally expensive (${\sim}$12 hours for the full
FNSPID corpus on a single GPU).  The \texttt{SentimentExtractor} class supports
a \texttt{cache\_key} parameter: when set, the extracted sentiment columns
(\texttt{sentiment\_label}, \texttt{sentiment\_score}, all three class
probabilities, and \texttt{sentiment\_value}) are serialised to a pickle file
under \texttt{data/cache/sentiment\_\{key\}.pkl}.  On subsequent calls, if the
cache file exists, the extraction step is skipped entirely and the cached
DataFrame is loaded and joined back to the input.  Additionally, per-ticker
Parquet files are saved under \texttt{data/cache/sentiment/\{TICKER\}\_sentiment.parquet}
to enable ticker-level selective recomputation.

\paragraph{Layer 3 --- Feature Cache (\texttt{FeatureCache}).}
Once technical indicators (V7) are computed and features are selected, the
resulting $(X, y, \text{feature\_cols})$ arrays are cached as pickle files under
\texttt{data/feature\_cache/}.  Cache files are keyed by
\texttt{(n\_stocks, start\_date, end\_date, version\_suffix)}.  This allows
experiments to skip the entire indicator computation and feature selection
pipeline, loading pre-built NumPy arrays directly.  This layer is especially
valuable given that indicator computation for 150 stocks over 8 years involves
${\sim}$300{,}000 rows $\times$ 100 indicators, which takes several minutes.

\begin{figure}[htbp]
\centering
\begin{verbatim}
data/
├── price_cache/
│   └── prices_2012-01-01_2019-12-31.pkl     ← Layer 1
├── cache/
│   ├── sentiment/
│   │   ├── AAPL_sentiment.parquet            ← Layer 2
│   │   ├── MSFT_sentiment.parquet
│   │   └── ...
│   └── sentiment_{key}.pkl                   ← Layer 2 (batch)
└── feature_cache/
    └── features_150stocks_2012_2019_v7.pkl   ← Layer 3
\end{verbatim}
\caption{Three-layer caching hierarchy.  Each layer eliminates a progressively
more expensive computation: network I/O (Layer~1), GPU-bound FinBERT inference
(Layer~2), and CPU-bound indicator computation with feature selection (Layer~3).}
\label{fig:cache-hierarchy}
\end{figure}

% ·····································
\subsection{Ticker Universe Curation}
\label{sec:ticker-curation}
% ·····································

Not all 200 candidate tickers provide reliable sentiment data.  A density
analysis pipeline computes each ticker's \textbf{sentiment density} --- the
fraction of business days within the training period that have at least one
non-neutral headline.  Tickers are ranked by density and the top~150 are
retained.

Additionally, a manual exclusion list of 90 tickers was compiled during the
density analysis phase.  These tickers exhibited one or more of the following
pathologies: delisted during the study period (e.g.\ YHOO, TWX), dominated
by ETF-level news rather than company-specific news (e.g.\ VNQ, XLK, SLV),
ticker collision with non-equity instruments (e.g.\ DGAZ, UGAZ), or
consistently noisy/irrelevant sentiment scores.

% ·····································
\subsection{Correlation Analysis: Sentiment vs.\ Returns}
\label{sec:correlation-analysis}
% ·····································

Before committing to end-to-end model training, we conducted a preliminary
correlation study to assess whether the raw sentiment signal carries any
predictive information.

For each ticker in our universe, we computed two Pearson correlation
coefficients:

\begin{enumerate}
  \item \textbf{Same-day correlation}: $\rho(\texttt{sent\_mean}_t,\; R_t)$,
        where $R_t = (P_t - P_{t-1}) / P_{t-1}$ is the daily return.  This
        captures whether sentiment and price move together contemporaneously.

  \item \textbf{Next-day (predictive) correlation}:
        $\rho(\texttt{sent\_mean}_t,\; R_{t+1})$, where $R_{t+1}$ is the
        following day's return.  This is the operationally relevant quantity,
        since a trading strategy can only act on today's sentiment to predict
        tomorrow's return.
\end{enumerate}

Neutral headlines were filtered out prior to correlation computation to avoid
diluting the signal.  Scatter plots with linear trend lines were generated for
each ticker for both same-day and next-day correlations, saved under
\texttt{reports/top40\_scatters/} and \texttt{reports/top40\_scatters\_next\_day/}
respectively.

\paragraph{Findings.}
The correlation analysis revealed that raw sentiment--return correlations are
\textbf{weak} across the board: same-day correlations typically fall in the range
$[-0.05, +0.10]$, and next-day correlations are even weaker, clustering near
zero.  This result is not unexpected --- the Efficient Market Hypothesis predicts
that publicly available sentiment should be rapidly priced in --- but it has an
important methodological implication: \emph{sentiment features cannot be expected
to act as standalone predictors}.  Their value, if any, must arise from
\textbf{nonlinear interactions} with technical indicators, which motivates the
neural-network approach rather than a simple linear combination.

% --------------------------------------------------------------------------
\section{Structural Model}
\label{sec:structural-model}
% --------------------------------------------------------------------------

This section describes the neural network architectures explored in this
project.  The progression from simple to complex architectures reflects an
empirical search for the right inductive bias for sentiment-augmented stock
prediction.

% ·····································
\subsection{Monolithic Architectures (Sentiment as a Regular Feature)}
\label{sec:monolithic}
% ·····································

The first approach to integrating sentiment treated it as ``just another
feature'': the sentiment columns were concatenated alongside the technical
indicators into a single $d$-dimensional input vector $\mathbf{x}_t \in
\mathbb{R}^d$ (where $d = 20 + 2 = 22$), and the entire vector was fed into a
standard recurrent network.

\subsubsection{SimpleTrendLSTM}

The simplest architecture uses a single-layer LSTM with hidden size 32:

\begin{equation}
  \mathbf{h}_t = \text{LSTM}(\mathbf{x}_t, \mathbf{h}_{t-1})
  \label{eq:simple-lstm}
\end{equation}

The final hidden state $\mathbf{h}_T$ is projected through a two-layer MLP
($32 \to 32 \to 3$) to produce trend logits.  Total parameters: ${\sim}$10K.

\subsubsection{AttentionLSTM}

Building on the SimpleTrendLSTM, the AttentionLSTM adds a multi-head
self-attention mechanism over the LSTM output sequence:

\begin{align}
  \mathbf{H} &= \text{LSTM}(\mathbf{X}) \in \mathbb{R}^{T \times h} \\
  \mathbf{Q} = \mathbf{H}\mathbf{W}_Q,\quad
  \mathbf{K} &= \mathbf{H}\mathbf{W}_K,\quad
  \mathbf{V} = \mathbf{H}\mathbf{W}_V \\
  \text{Attn}(\mathbf{Q}, \mathbf{K}, \mathbf{V}) &=
    \text{softmax}\!\left(\frac{\mathbf{Q}\mathbf{K}^\top}{\sqrt{d_k}}\right)
    \mathbf{V}
  \label{eq:attention}
\end{align}

The attention module (\texttt{SentimentAttention}, 4 heads by default) learns to
weight different timesteps in the lookback window, allowing the model to focus
on the most informative days rather than relying solely on the final hidden
state.  Total parameters: ${\sim}$50K.

\paragraph{Limitation.}
In both monolithic architectures, the LSTM's input gate, forget gate, and cell
state treat every input dimension identically.  This means the two sentiment
features (which are sparse, noisy, and operate on a fundamentally different
timescale than technical indicators) compete with 20 well-behaved technical
features for the same recurrent capacity.  Experimentally, this led to the
sentiment signal being ``drowned out'' during gradient updates.

% ·····································
\subsection{Dual-Branch Hybrid Architecture}
\label{sec:hybrid}
% ·····································

To address the monolithic architecture's inability to treat heterogeneous
feature types appropriately, we developed a \textbf{dual-branch hybrid
architecture} that processes technical indicators and sentiment features through
separate neural pathways before fusing them.

\subsubsection{Architecture Overview}

The model accepts a single concatenated input tensor $\mathbf{x} \in
\mathbb{R}^{B \times T \times (n_{\text{tech}} + n_{\text{sent}})}$ and splits
it internally by feature index:

\begin{equation}
  \mathbf{x}_{\text{tech}} = \mathbf{x}[:, :, :n_{\text{tech}}],\qquad
  \mathbf{x}_{\text{sent}} = \mathbf{x}[:, :, n_{\text{tech}}:]
\end{equation}

\noindent
This internal-split design preserves full compatibility with the existing data
pipeline, scaler, and walk-forward loop --- no upstream changes are required.

\subsubsection{Branch 1: Technical Branch (LSTM + Multi-Head Attention)}

The technical branch processes the 20 SHAP-selected indicators:

\begin{align}
  \hat{\mathbf{x}}_{\text{tech}} &= \text{LayerNorm}(\mathbf{x}_{\text{tech}}) \\
  \mathbf{H}_{\text{tech}} &= \text{LSTM}_{2\text{-layer, }h{=}128}
    (\hat{\mathbf{x}}_{\text{tech}})
    \in \mathbb{R}^{B \times T \times 128} \\
  \mathbf{A}_{\text{tech}} &= \text{MultiHeadAttn}_{4\text{ heads}}
    (\mathbf{H}_{\text{tech}})
    \in \mathbb{R}^{B \times T \times 128} \\
  \mathbf{z}_{\text{tech}} &= \mathbf{A}_{\text{tech}}[:, -1, :]
    \in \mathbb{R}^{B \times 128}
\end{align}

A two-layer LSTM with hidden size 128 captures temporal dependencies in the
technical indicators.  The 4-head self-attention mechanism applied to the LSTM
outputs learns which timesteps within the 20-day lookback window are most
informative for the prediction task.  The attended representation at the last
timestep is extracted as the branch output.

\subsubsection{Branch 2: Sentiment Branch (Lightweight GRU)}

The sentiment branch processes the 2 RF-selected sentiment features:

\begin{align}
  \hat{\mathbf{x}}_{\text{sent}} &= \text{LayerNorm}(\mathbf{x}_{\text{sent}}) \\
  \mathbf{H}_{\text{sent}},\; \mathbf{h}_{\text{sent}} &=
    \text{GRU}_{1\text{-layer, }h{=}32}(\hat{\mathbf{x}}_{\text{sent}}) \\
  \mathbf{z}_{\text{sent}} &= \mathbf{h}_{\text{sent}}[-1]
    \in \mathbb{R}^{B \times 32}
\end{align}

Several deliberate design choices distinguish this branch:

\begin{itemize}
  \item \textbf{GRU instead of LSTM.}  A GRU (Gated Recurrent Unit) has
        approximately $25\%$ fewer parameters than an LSTM of the same hidden
        size because it lacks a separate cell state ($\mathbf{c}_t$).  With only
        2 input features, an LSTM's additional capacity offers no benefit and
        increases the risk of overfitting.

  \item \textbf{Small hidden size ($h = 32$ vs.\ $h = 128$).}  The sentiment
        branch is intentionally constrained to ${\sim}$3{,}000 parameters so
        that it cannot dominate gradient updates.  Sentiment should act as a
        \emph{learned modifier} of the technical signal, not an equal partner.

  \item \textbf{No attention.}  The sentiment features (\texttt{sent\_surprise},
        \texttt{sent\_ema\_20d}) are already derived summary statistics that
        aggregate information over 20-day windows.  Adding attention on top of
        these would introduce unnecessary complexity.

  \item \textbf{Single layer.}  Deeper recurrence is unnecessary for 2-dimensional
        input and would further increase overfitting risk.
\end{itemize}

\subsubsection{Fusion Layer}

The outputs of both branches are concatenated and processed through a shared
dense network:

\begin{align}
  \mathbf{z}_{\text{fused}} &= [\mathbf{z}_{\text{tech}} \;\|\;
    \mathbf{z}_{\text{sent}}]
    \in \mathbb{R}^{B \times 160} \\
  \mathbf{f} &= \text{Dropout}_{0.3}\bigl(\text{ReLU}\bigl(
    \mathbf{W}_{\text{fc}}\, \mathbf{z}_{\text{fused}} + \mathbf{b}_{\text{fc}}
    \bigr)\bigr)
    \in \mathbb{R}^{B \times 64}
\end{align}

The fusion layer learns to weight the relative importance of technical and
sentiment representations adaptively.

\subsubsection{Task-Specific Output Heads}

\paragraph{Classification Head (HybridAttentionLSTM).}
For the 3-class trend prediction task:

\begin{equation}
  \hat{\mathbf{y}} = \mathbf{W}_2\, \text{ReLU}(\mathbf{W}_1\, \mathbf{f} +
    \mathbf{b}_1) + \mathbf{b}_2
    \in \mathbb{R}^{B \times 3}
    \label{eq:cls-head}
\end{equation}

with an intermediate projection $64 \to 32 \to 3$ and dropout of $0.15$ between
layers.

\paragraph{Regression Head (HybridRegressionLSTM).}
For the vol-adjusted forward return prediction task:

\begin{equation}
  \hat{y} = \mathbf{w}^\top \text{ReLU}(\mathbf{W}_1\, \mathbf{f} +
    \mathbf{b}_1) + b
    \in \mathbb{R}^{B}
    \label{eq:reg-head}
\end{equation}

with the same $64 \to 32 \to 1$ projection, producing a single scalar output.

\subsubsection{Graceful Degradation}

When the \texttt{--no-sentiment} flag is passed, the model is instantiated with
$n_{\text{sent}} = 0$.  In this configuration, the sentiment branch is
automatically disabled (via a \texttt{has\_sentiment\_branch} flag) and the
fusion layer input size reduces from 160 to 128.  This enables controlled A/B
experiments where the \emph{only} difference is the presence or absence of the
sentiment branch.

% ·····································
\subsection{Weight Initialisation}
\label{sec:weight-init}
% ·····································

Proper weight initialisation is critical for training stability in recurrent
networks.  We apply different initialisation strategies to different parameter
types, reflecting the distinct roles they play:

\paragraph{Xavier Uniform Initialisation (Input-to-Hidden Weights).}
All \texttt{weight\_ih} (input-to-hidden) parameters in both the LSTM and GRU
layers are initialised using Xavier (Glorot) uniform
initialisation~\cite{glorot2010understanding}:

\begin{equation}
  W_{ij} \sim \mathcal{U}\!\left(-\sqrt{\frac{6}{n_{\text{in}} + n_{\text{out}}}},\;
  +\sqrt{\frac{6}{n_{\text{in}} + n_{\text{out}}}}\right)
\end{equation}

Xavier initialisation sets the variance of weights such that the variance of
activations is preserved across layers during the forward pass.  This prevents
the signal from either exploding or vanishing at initialisation, which is
especially important for the first layer of a recurrent network where the input
distribution is known to be standardised.

\paragraph{Orthogonal Initialisation (Hidden-to-Hidden Weights).}
All \texttt{weight\_hh} (hidden-to-hidden, recurrent) parameters are initialised
as orthogonal matrices~\cite{saxe2014exact}:

\begin{equation}
  \mathbf{W}_{\text{hh}} = \mathbf{U}\mathbf{V}^\top, \quad
  \text{where } \mathbf{U}\boldsymbol{\Sigma}\mathbf{V}^\top =
  \text{SVD}(\mathbf{A}),\; A_{ij} \sim \mathcal{N}(0, 1)
\end{equation}

Orthogonal matrices have all singular values equal to 1, meaning that repeated
matrix--vector multiplication (as occurs across time steps in an RNN)
\emph{preserves the norm} of the hidden state vector.  This directly combats the
vanishing/exploding gradient problem that plagues recurrent networks: without
orthogonal initialisation, the spectral norm of $\mathbf{W}_{\text{hh}}$
deviating from 1 causes exponential amplification or attenuation over the
20-step lookback window.

\paragraph{Zero-Initialised Biases.}
All bias terms are initialised to zero.  Combined with the above weight
initialisations, this ensures that at the start of training the gates in the
LSTM (input, forget, output) and GRU (reset, update) are approximately
balanced, avoiding pathological initial states where the forget gate is fully
open or fully closed.

\paragraph{Dense Layers.}
All \texttt{nn.Linear} layers in the fusion network and output heads use Xavier
uniform initialisation with zero biases, consistent with standard practice for
feed-forward networks with ReLU activations.

% ·····································
\subsection{Loss Functions}
\label{sec:loss-functions}
% ·····································

The choice of loss function was guided by the specific challenges of each task
formulation.

\paragraph{Classification: Cross-Entropy with Label Smoothing.}
For the 3-class trend prediction task, we use \textbf{Cross-Entropy Loss} with
a label smoothing factor of $\epsilon = 0.1$:

\begin{equation}
  \mathcal{L}_{\text{CE}} = -\sum_{c=1}^{3}
    \tilde{y}_c \log \hat{y}_c, \qquad
  \tilde{y}_c = (1 - \epsilon)\, y_c + \frac{\epsilon}{3}
  \label{eq:ce-loss}
\end{equation}

Cross-Entropy is the standard loss for multi-class classification because it
directly optimises the log-likelihood of the correct class.  We chose it over
alternatives such as Focal Loss because our class distribution (down/neutral/up)
is roughly balanced by design (the return thresholds $\pm 0.5\%$ are calibrated
to produce near-equal class frequencies).

Label smoothing is applied because financial trend labels are inherently noisy:
a daily return of $+0.51\%$ (classified as ``up'') and $+0.49\%$ (classified as
``neutral'') are essentially the same event.  By softening the target
distribution, label smoothing discourages the model from becoming overconfident
about these boundary cases and improves generalisation.

\paragraph{Regression: Huber Loss.}
For the vol-adjusted return prediction task, we use \textbf{Huber Loss} (also
known as Smooth L1 Loss) with $\delta = 1.0$:

\begin{equation}
  \mathcal{L}_{\text{Huber}}(y, \hat{y}) =
  \begin{cases}
    \frac{1}{2}(y - \hat{y})^2 & \text{if } |y - \hat{y}| \leq \delta \\
    \delta\, |y - \hat{y}| - \frac{1}{2}\delta^2 & \text{otherwise}
  \end{cases}
  \label{eq:huber-loss}
\end{equation}

We chose Huber Loss over standard Mean Squared Error (MSE) because financial
returns are heavy-tailed: extreme returns (e.g.\ earnings surprises, flash
crashes) produce large residuals that, under MSE, generate disproportionately
large gradients.  Huber Loss transitions from quadratic (L2) behaviour near
zero to linear (L1) behaviour for large residuals, effectively clipping the
influence of outliers without discarding them entirely.  This yields more
stable training dynamics compared to MSE, which we observed to produce
oscillating loss curves in early experiments.

We also explored \textbf{MSE Loss} and \textbf{Log-Cosh Loss} in earlier
iterations.  MSE was abandoned due to sensitivity to outlier returns.
Log-Cosh --- which is approximately $\frac{1}{2}x^2$ for small $x$ and
$|x| - \log 2$ for large $x$ --- performed comparably to Huber but offered
no clear advantage while being less interpretable.

% --------------------------------------------------------------------------
\section{Dynamic Model}
\label{sec:dynamic-model}
% --------------------------------------------------------------------------

The dynamic model describes how the structural model is trained, evaluated,
and evolved over time.

% ·····································
\subsection{Walk-Forward (Rolling-Origin) Validation}
\label{sec:walk-forward}
% ·····································

Standard $k$-fold cross-validation is inappropriate for time-series data because
it violates the temporal ordering of observations: a model trained on 2019 data
and tested on 2017 data would exploit information from the future.  We therefore
employ \textbf{walk-forward validation} with an \textbf{expanding window}
protocol:

\begin{enumerate}
  \item The date range 2017-Q1 through 2019-Q4 is divided into 12 quarterly
        folds.  The fold boundaries are:
        $\{$2017-01, 2017-04, 2017-07, 2017-10, 2018-01, \ldots, 2019-10$\}$.

  \item For fold $k$ with validation period $[t_k^{\text{start}},\,
        t_k^{\text{end}})$:
        \begin{itemize}
          \item \textbf{Training set}: all data from 2012-01-01 up to
                $t_k^{\text{start}}$ (expanding window).
          \item \textbf{Validation set}: data from
                $[t_k^{\text{start}},\, t_k^{\text{end}})$.
        \end{itemize}

  \item A \textbf{fresh model} is initialised and trained from scratch for each
        fold.  There is no warm-starting or weight transfer between folds.  This
        prevents information leakage from future folds and simulates the
        realistic deployment scenario where a model is periodically retrained
        from scratch.

  \item Each fold's scaler is fitted exclusively on that fold's training data.
\end{enumerate}

\noindent
After all 12 folds, the per-fold predictions are concatenated and evaluated as a
single continuous prediction stream, yielding a final aggregate metric that
represents realistic out-of-sample performance.

\begin{figure}[htbp]
\centering
\begin{verbatim}
Fold  Training Data (expanding)    Validation Quarter
─── ─────────────────────────── ──────────────────────
  1  2012-01 ──────────── 2017-01  [2017-Q1]
  2  2012-01 ──────────── 2017-04  [2017-Q2]
  3  2012-01 ──────────── 2017-07  [2017-Q3]
  4  2012-01 ──────────── 2017-10  [2017-Q4]
  5  2012-01 ──────────── 2018-01  [2018-Q1]
  6  2012-01 ──────────── 2018-04  [2018-Q2]
  7  2012-01 ──────────── 2018-07  [2018-Q3]
  8  2012-01 ──────────── 2018-10  [2018-Q4]
  9  2012-01 ──────────── 2019-01  [2019-Q1]
 10  2012-01 ──────────── 2019-04  [2019-Q2]
 11  2012-01 ──────────── 2019-07  [2019-Q3]
 12  2012-01 ──────────── 2019-10  [2019-Q4]
\end{verbatim}
\caption{Walk-forward validation schedule (12 quarterly folds).  Each fold's
training window expands by one quarter while the validation quarter advances
forward.}
\label{fig:walk-forward}
\end{figure}

% ·····································
\subsection{Sequence Construction}
\label{sec:sequences}
% ·····································

From each stock's time series, fixed-length input sequences are constructed via a
sliding window:

\begin{equation}
  \mathbf{X}_i = [\mathbf{x}_{i},\; \mathbf{x}_{i+1},\; \ldots,\;
    \mathbf{x}_{i+L-1}] \in \mathbb{R}^{L \times d}, \qquad
  y_i = \text{target}_{i+L}
\end{equation}

where $L = 20$ trading days (approximately one calendar month) is the lookback
window and $d = 22$ is the feature dimensionality.  Sequences are assigned to
the fold whose validation period contains the \emph{target date} $t_{i+L}$.
The lookback window naturally extends before the fold boundary, ensuring that
the model has sufficient context.

Multi-stock data is pooled: sequences from all training-set tickers are
concatenated into a single array, shuffled, and batched.  This cross-sectional
pooling forces the model to learn \emph{universal} patterns rather than
stock-specific idiosyncrasies.

% ·····································
\subsection{Training Protocol}
\label{sec:training-protocol}
% ·····································

Each fold trains independently with the following configuration:

\begin{table}[htbp]
\centering
\caption{Training hyperparameters.}
\label{tab:training-config}
\begin{tabular}{@{}ll@{}}
\toprule
\textbf{Parameter}           & \textbf{Value} \\
\midrule
Optimiser                    & AdamW ($\beta_1{=}0.9$, $\beta_2{=}0.999$) \\
Learning rate                & $5 \times 10^{-4}$ \\
Weight decay                 & $1 \times 10^{-4}$ \\
Batch size                   & 1024 (daily), 512 (weekly) \\
Max epochs                   & 100 \\
Early stopping patience      & 15 epochs \\
LR scheduler                 & ReduceLROnPlateau (factor 0.5, patience 5) \\
Minimum LR                   & $1 \times 10^{-6}$ \\
Gradient clipping             & Max norm 1.0 \\
Mixed precision              & FP16 via \texttt{torch.amp.GradScaler} \\
Dropout                      & 0.3 (recurrent), 0.15 (output heads) \\
Label smoothing (cls.\ only) & 0.1 \\
\bottomrule
\end{tabular}
\end{table}

\paragraph{Mixed-Precision Training.}
All GPU-accelerated folds use \texttt{torch.amp.autocast} (FP16) with
\texttt{GradScaler} to halve memory consumption and increase throughput.  The
\texttt{GradScaler} dynamically scales the loss to prevent underflow in FP16
gradients, and \texttt{unscale\_} is called before gradient clipping to ensure
the clip threshold operates on the true gradient magnitudes.

\paragraph{Early Stopping.}
The validation loss is monitored after each epoch.  If it fails to improve for
15 consecutive epochs, training is halted and the best model checkpoint (lowest
validation loss) is restored.  This prevents overfitting on the training set
while avoiding premature termination.

\paragraph{Normalisation.}
A \texttt{StandardScaler} (z-score normalisation) is fitted exclusively on each
fold's training data and applied to both training and validation splits.  This
ensures zero mean and unit variance for all features, which is critical for
LSTM stability since the gates' sigmoid activations are sensitive to input
scale.

% ·····································
\subsection{A/B Experiment Design}
\label{sec:ab-design}
% ·····································

To isolate the marginal contribution of sentiment, every experiment is run in
two configurations:

\begin{enumerate}
  \item \textbf{Sentiment ON}: the full 22-feature vector (20 technical + 2
        sentiment) is used, and the dual-branch architecture processes them
        through separate pathways.

  \item \textbf{Sentiment OFF}: the same ticker universe, date range, and
        hyperparameters are used, but the 2 sentiment features are zeroed out.
        The model is instantiated with $n_{\text{sent}} = 0$, which disables the
        sentiment branch entirely.
\end{enumerate}

\noindent
By holding the ticker universe constant (even in the no-sentiment condition),
we ensure that any performance difference is attributable solely to the
sentiment signal, not to a different stock selection.

Additionally, the monolithic architecture (AttentionLSTM with sentiment as
regular features) is run as a third configuration to isolate the contribution of
the architectural change (dual-branch) from the contribution of the data
(sentiment features).


% ============================================================================
\chapter{Experimentation Environment and Experiment Design}
\label{ch:experiments}
% ============================================================================

% ·····································
\section{Hardware and Software Environment}
\label{sec:hardware}
% ·····································

All experiments were executed on the following configuration:

\begin{table}[htbp]
\centering
\caption{Hardware and software environment.}
\label{tab:environment}
\begin{tabular}{@{}ll@{}}
\toprule
\textbf{Component}  & \textbf{Specification} \\
\midrule
GPU                  & NVIDIA GeForce RTX 2060 (6\,GB VRAM) \\
CPU                  & Intel / AMD (host) \\
RAM                  & 16\,GB DDR4 \\
OS                   & Windows \\
Python               & 3.10+ \\
PyTorch              & 2.x with CUDA \\
Key Libraries        & \texttt{transformers} (FinBERT), \texttt{pandas\_ta},
                       \texttt{scikit-learn}, \texttt{shap}, \texttt{scipy} \\
\bottomrule
\end{tabular}
\end{table}

\noindent
The 6\,GB VRAM constraint directly informed several design decisions:
FinBERT inference batch size was limited to 16, LSTM hidden sizes were kept at
128 (rather than the 256--512 common in NLP), and mixed-precision training
(FP16) was mandatory for the larger walk-forward experiments.

% ·····································
\section{Experiment Progression}
\label{sec:experiment-progression}
% ·····································

The experimental campaign comprised six phases, each informed by the failures
and insights of the preceding phase.  Table~\ref{tab:experiment-phases}
summarises the progression.

\begin{table}[htbp]
\centering
\caption{Six-phase experiment progression.}
\label{tab:experiment-phases}
\begin{tabular}{@{}clp{8cm}@{}}
\toprule
\textbf{Phase} & \textbf{Name} & \textbf{Description} \\
\midrule
1 & Single-Stock Baseline   & LSTM trained per-stock on OHLCV; established
                                that per-stock models overfit badly. \\
2 & Pooled Baseline (V7--V10) & Multi-stock pooled training with
                                ComprehensiveIndicatorsV7; SHAP feature
                                selection; achieved stable but modest
                                baselines. \\
3 & Sentiment Correlation   & Pearson correlation of FinBERT sentiment vs.\
                                same-day and next-day returns; confirmed
                                weak linear signal. \\
4 & Monolithic Sentiment    & Sentiment features concatenated as additional
                                LSTM inputs; walk-forward validation; no
                                improvement over no-sentiment baseline. \\
5 & Dual-Branch Hybrid      & Separate LSTM+Attention for technicals and GRU
                                for sentiment; walk-forward validation;
                                marginal improvements in some folds. \\
6 & Architecture Comparison & Head-to-head A/B: monolithic vs.\ dual-branch
                                vs.\ no-sentiment; established that
                                architectural separation slightly outperforms
                                monolithic but sentiment overall does not
                                produce statistically significant
                                improvement. \\
\bottomrule
\end{tabular}
\end{table}

% ·····································
\section{Classification Experiment Design}
\label{sec:cls-experiment}
% ·····································

\subsection{Task Definition}

The classification task predicts a 3-class trend label for each stock on each
trading day:

\begin{equation}
  y_t =
  \begin{cases}
    0 \;(\text{Down})    & \text{if } R_{t+1} < -0.5\% \\
    1 \;(\text{Neutral}) & \text{if } -0.5\% \leq R_{t+1} \leq +0.5\% \\
    2 \;(\text{Up})      & \text{if } R_{t+1} > +0.5\%
  \end{cases}
  \label{eq:trend-label}
\end{equation}

where $R_{t+1} = (P_{t+1} - P_t) / P_t$ is the next-day simple return.

\subsection{Metrics}

\begin{itemize}
  \item \textbf{Accuracy}: fraction of correctly classified samples.
  \item \textbf{Macro F1-Score}: unweighted mean of per-class F1 scores,
        ensuring that minority classes receive equal importance.
  \item \textbf{Per-class Precision and Recall}: to diagnose whether the model
        is biased toward the neutral class (a common failure mode).
  \item \textbf{Cohen's Kappa}: agreement beyond chance, since a naive
        always-neutral classifier would achieve ${\sim}33\%$ accuracy.
\end{itemize}

\subsection{Sentiment Features (Classification)}

The two features selected by RF permutation importance for the classification
task are:

\begin{enumerate}
  \item \texttt{sent\_surprise}: $\texttt{sent\_mean}_t - \text{MA}_{20}(\texttt{sent\_mean})$.
        Measures deviation from the 20-day sentiment baseline.  Captures days
        where sentiment is abnormally positive or negative relative to the
        recent regime.

  \item \texttt{sent\_ema\_20d}: $\text{EMA}_{20}(\texttt{sent\_mean})$.
        Captures the long-horizon sentiment regime, smoothing out daily
        fluctuations.
\end{enumerate}

% ·····································
\section{Regression Experiment Design}
\label{sec:reg-experiment}
% ·····································

\subsection{Task Definition}

The regression task predicts the \textbf{volatility-adjusted forward return}:

\begin{equation}
  y_t = \frac{R_{t+1}}{\sigma_{t}^{(20)}}, \qquad
  \sigma_{t}^{(20)} = \text{std}(R_{t-19}, \ldots, R_t)
  \label{eq:vol-adj-return}
\end{equation}

where $\sigma_{t}^{(20)}$ is the trailing 20-day return volatility.  This
normalisation makes the target comparable across stocks with different
volatility profiles and across different market regimes.  The target is clipped
to $[-10, +10]$ to prevent extreme outliers from destabilising training.

\subsection{Metrics}

\begin{itemize}
  \item \textbf{Mean Squared Error (MSE)} and \textbf{Mean Absolute Error (MAE)}:
        standard regression metrics.
  \item \textbf{$R^2$ (coefficient of determination)}: fraction of variance
        explained by the model.  $R^2 > 0$ indicates improvement over predicting
        the mean.
  \item \textbf{Information Coefficient (IC)}: Spearman rank correlation between
        predicted and actual returns across all stocks at each time point.  A
        positive IC indicates that the model's cross-sectional ranking of stocks
        has predictive value, even if the point predictions are inaccurate.
  \item \textbf{Directional Accuracy}: fraction of days where the model correctly
        predicts the sign of the return.
\end{itemize}

\subsection{Sentiment Features (Regression)}

The regression task uses a different pair of sentiment features, also selected
by RF permutation importance but with $R^2$ scoring:

\begin{enumerate}
  \item \texttt{sent\_count\_log}: $\log(1 + \text{news count})$.  Captures
        abnormal news activity, which often precedes large price moves
        regardless of sentiment direction.

  \item \texttt{sent\_strength}: $|\texttt{sent\_mean}| \times
        \log(1 + \text{count})$.  An interaction term that combines sentiment
        polarity magnitude with news volume, capturing the intuition that
        a strongly negative day \emph{with many articles} is more informative
        than a single negative headline.
\end{enumerate}

% ·····································
\section{Controlled Comparison Design}
\label{sec:controlled}
% ·····································

Each experiment configuration is run in a controlled factorial design:

\begin{table}[htbp]
\centering
\caption{Factorial experiment design.  Each row represents a distinct
experimental run, all using the same ticker universe, date range, and
hyperparameters.}
\label{tab:factorial}
\begin{tabular}{@{}lllp{6cm}@{}}
\toprule
\textbf{Run} & \textbf{Architecture} & \textbf{Sentiment} & \textbf{Purpose} \\
\midrule
A & HybridAttentionLSTM  & ON  & Full proposed system \\
B & HybridAttentionLSTM  & OFF & Isolates sentiment contribution \\
C & AttentionLSTM (mono) & ON  & Isolates architectural contribution \\
D & AttentionLSTM (mono) & OFF & Baseline (no sentiment, no hybrid) \\
\bottomrule
\end{tabular}
\end{table}

\noindent
Comparing A vs.\ B reveals the \emph{marginal value of sentiment} given the
hybrid architecture.  Comparing A vs.\ C reveals the \emph{marginal value of
the hybrid architecture} given sentiment.  Comparing A vs.\ D reveals the
\emph{combined effect} of both innovations.

% ·····································
\section{Auxiliary Analyses}
\label{sec:auxiliary}
% ·····································

Beyond the main experiments, several auxiliary analyses were conducted to
deepen understanding of the system:

\paragraph{SHAP Importance Analysis.}
SHAP values were computed for the technical feature selection model (XGBoost)
to produce a ranked list of the most informative indicators.  The top 20
features were retained; the remaining ${\sim}$80 were discarded.

\paragraph{Sentiment Density Analysis.}
A ticker-by-ticker density analysis computed the fraction of business days with
at least one non-neutral headline.  This analysis produced the exclusion list
of 90 noisy tickers and the density-based ranking used for ticker selection.

\paragraph{Sentiment--Return Correlation Analysis.}
As described in Section~\ref{sec:correlation-analysis}, Pearson correlations
were computed between daily sentiment and both same-day and next-day returns
for all tickers, with scatter-plot visualisations produced for each.

\paragraph{Hyperparameter Sensitivity.}
Learning rate, hidden size, dropout, and sequence length were varied in a grid
search during the pooled baseline phase.  The final hyperparameters
(Table~\ref{tab:training-config}) represent the configuration that achieved
the lowest validation loss across the majority of folds.

\paragraph{Model Size Ablation.}
The sentiment branch hidden size was varied (16, 32, 64) to verify that the
small branch (32) achieves the best trade-off between capacity and overfitting.


% ============================================================================
% Placeholder for bibliography entries referenced above
% ============================================================================
% \begin{thebibliography}{99}
% \bibitem{fnspid2024} ...
% \bibitem{finbert2019} ...
% \bibitem{devlin2019bert} ...
% \bibitem{glorot2010understanding} ...
% \bibitem{saxe2014exact} ...
% \bibitem{lundberg2017shap} ...
% \end{thebibliography}
